\section*{Chapitre \ref{ch:qm:linear-models-under-stress} --- Modèles linéaires sous contrainte}

\begin{solution}{exo:qm:linear-models-under-stress:1}
$R_j^2 = 0.81$, donc $\mathrm{VIF} = 1/0.19 = 5.3$ et chaque \omterm{def:qm:estimation:estimator}{erreur type} croît d'un facteur $\sqrt{5.3} = 2.3$.
\end{solution}

\begin{solution}{exo:qm:linear-models-under-stress:2}
$\lambda = 1/(1 + 0.25) = 0.8$, donc les \omterm{def:qm:linear-models-under-stress:ols}{moindres carrés ordinaires} donnent $0.8 \times 2 = 1.6$.
\end{solution}

\begin{solution}{exo:qm:linear-models-under-stress:3}
$\sum_ih_{ii} = \operatorname{tr}H = \operatorname{tr}\bigl(X(X^\top X)^{-1}X^\top\bigr) = \operatorname{tr}\bigl((X^\top X)^{-1}X^\top X\bigr) = \operatorname{tr}I_k = k$.
\end{solution}

\begin{solution}{exo:qm:linear-models-under-stress:4}
Avec $X^\top X = nI$, la fonction objectif se sépare~: $\frac1{2n}\lvert y - X\beta\rvert^2 = \mathrm{const} + \frac12\sum_j(\beta_j - b_j)^2$, où $b = X^\top y/n$ est le
coefficient des moindres carrés. Chaque coordonnée minimise $\frac12(\beta_j - b_j)^2 + \lambda|\beta_j|$, dont la solution est $\operatorname{sign}(b_j)\max(|b_j| -
\lambda, 0)$~: seuillage doux.
\end{solution}

\begin{solution}{exo:qm:linear-models-under-stress:5}
La variance de l'erreur de niveau vaut $-\gamma(1) = 0.0196$ (écart type 0.14), donc $\hat\lambda = (0.199 - 2 \times 0.0196)/0.199 = 0.80$. Au pas hebdomadaire, la vraie
variance vaut $5 \times (0.199 - 0.0392) = 0.80$, à comparer aux mêmes $0.0392$ d'erreur~: $\lambda_5 = 0.80/0.839 = 0.95$.
\end{solution}

\begin{solution}{exo:qm:linear-models-under-stress:6}
Avec $X^\top X \approx n\begin{psmallmatrix}1 & \rho\\ \rho & 1\end{psmallmatrix}$, $\Var(\hat\beta_1) = \sigma^2/(n(1 - \rho^2)) = 16.9\sigma^2/n$, tandis que
$\Var(\hat\beta_1 + \hat\beta_2) = 2\sigma^2/(n(1 + \rho)) = 1.02\sigma^2/n$. La somme, c'est-à-dire l'exposition à la direction commune, est déterminée environ quatre fois
plus précisément (en écart type) que chacun des deux bêtas~; la différence, de variance $2\sigma^2/(n(1 - \rho)) = 66.7\sigma^2/n$, est l'endroit où
se loge le bruit.
\end{solution}

\begin{solution}{exo:qm:linear-models-under-stress:7}
\texttt{panel\_truth()}~: écart type vrai 0.105~; classique 0.026 et White 0.026~; \omterm{def:qm:linear-models-under-stress:fm}{erreurs types groupées} (clustered) par entreprise 0.101~; Fama--MacBeth 0.023 (0.022 avec trois
retards de Newey--West). L'effet entreprise est commun à la pente de chaque mois, si bien que ni la formule classique ni la dispersion temporelle des pentes mensuelles
ne le voit~; le regroupement par entreprise, lui, le capte.
\end{solution}

\begin{solution}{exo:qm:linear-models-under-stress:8}
Des plis aléatoires placent des jours postérieurs dans l'ensemble d'apprentissage de chaque jour de test, et comme les cibles de vingt jours se chevauchent sur dix-neuf jours,
des observations voisines partagent l'essentiel de leur cible~: le modèle est testé sur des rendements qu'il a en réalité déjà vus. Il faut utiliser des plis croissants (ou glissants) dans
l'ordre temporel, avec une purge d'au moins vingt jours entre l'apprentissage et le test, et s'attendre à un $R^2$ nettement plus faible~; compter aussi les pénalités et
les spécifications testées (chapitre~12).
\end{solution}

\begin{solution}{pb:qm:linear-models-under-stress:1}
\textbf{1.} 0.68.
\textbf{2.} $\lambda = 0.16/(0.16 + 2 \times 0.14^2) = 0.80$, donc en moyenne $0.80 \times 0.85 = 0.683$.
\textbf{3.} Moyenne 0.683, écart type 0.017.
\textbf{4.} Le bruit dans la variable explicative biaise la pente~; davantage de données réduisent la variance autour de la mauvaise valeur, pas le biais.
\textbf{5.} 0.086 points par jour contre 0.049 pour le ratio vrai, soit 74~\% de plus.
\textbf{6.} $-0.098$~; $\hat\lambda = 1 + 2 \times (-0.098) = 0.80$.
\textbf{7.} $0.68/0.80 = 0.85$, résidu 0.049.
\textbf{8.} $\lambda_5 = 0.95$~; le ratio hebdomadaire vaut 0.80, résidu 0.053.
\textbf{9.} Corrigé~: moyenne 0.862, écart type 0.095. Hebdomadaire~: moyenne 0.810, écart type 0.022.
\textbf{10.} 0.74~; elle suppose que les variances des erreurs sur l'obligation et sur le contrat à terme sont égales, alors qu'ici celle du contrat à terme vaut 0.039 et le résidu de l'obligation
est de 0.0025.
\textbf{11.} $0.16/(0.16 + 2 \times 0.07^2) = 0.94$.
\textbf{12.} Cela supprime le problème~: aucune erreur de niveau, aucune atténuation, et le ratio des \omterm{def:qm:linear-models-under-stress:ols}{moindres carrés ordinaires} quotidiens est sans biais.
\textbf{13.} Mensuellement, $\lambda_{21} = 21 \times 0.16/(21 \times 0.16 + 0.039) = 0.99$, presque aucun biais, mais seulement 24 observations sur deux ans, donc un
ratio bruité~; la régression hebdomadaire est le compromis.
\textbf{14.} Non~: la colinéarité traduit un manque de variation indépendante, et non un bruit dans une variable explicative. Agréger ne sépare pas les facteurs~; seul le fait de rapporter
leur somme, de pénaliser ou de trouver des données où ils évoluent séparément, y parvient.
\textbf{15.} Comparer au ratio fondé sur la duration, tester le P\&L couvert hors échantillon, et vérifier la stabilité du ratio selon les périodes et
les horizons.
\textbf{16.} Le ratio corrigé ou fondé sur la duration, proche de 0.85, recoupé par le 0.80 hebdomadaire~; et non le 0.68 des \omterm{def:qm:linear-models-under-stress:ols}{moindres carrés ordinaires} quotidiens.
\textbf{17.} Il faut enregistrer les cours médians synchrones du contrat à terme au moment de la valorisation de l'obligation.
\textbf{18.} Lorsque l'objectif est la prévision à partir de la variable explicative bruitée elle-même (la pente atténuée est alors le meilleur prédicteur linéaire), plutôt que la
relation structurelle.
\textbf{19.} Résultat nommé~: \emph{le ratio de couverture qui a rétréci}~: des cours du contrat à terme, périmés et bruités par le bid-ask, atténuent de 0.80 le ratio de couverture estimé par \omterm{def:qm:linear-models-under-stress:ols}{moindres carrés ordinaires} sur données quotidiennes, qui passe de
0.85 à 0.68~; le ratio corrigé à partir de la première autocorrélation propre au contrat à terme vaut 0.85~; le risque résiduel est de 0.086 point par jour à 0.68, contre 0.049
à 0.85 (0.053 au ratio hebdomadaire de 0.80).
\textbf{20.} Il biaise la pente vers zéro, alors que le bruit dans la variable expliquée ne fait qu'ajouter de la variance.
\end{solution}

\begin{solution}{iq:qm:linear-models-under-stress:1}
La variance de chaque coefficient est multipliée par $1/(1 - 0.97^2) = 17$~; les coefficients individuels oscillent et peuvent changer de signe, alors que leur somme est bien
déterminée. Rapporter la somme ou l'exposition en composantes principales, pénaliser (\omterm{def:qm:linear-models-under-stress:penalties}{régression ridge}), ou retirer un facteur.
\iqlookfor{Le \omterm{def:qm:linear-models-under-stress:vif}{facteur d'inflation de la variance}, la combinaison stable et un remède adapté à la question posée.}
\end{solution}

\begin{solution}{iq:qm:linear-models-under-stress:2}
Atténuation due au bruit de la variable explicative~: cours asynchrones ou périmés, rebond bid-ask dans les transactions du contrat à terme. Également une relation qui change (un
échantillon couvrant des régimes différents), ou un mauvais horizon. Vérifier la première autocorrélation du contrat à terme, relancer l'estimation sur les variations hebdomadaires, et
utiliser des cours synchrones.
\iqlookfor{Les \omterm{def:qm:linear-models-under-stress:eiv}{erreurs sur les variables} comme premier suspect, avec un diagnostic et une correction.}
\end{solution}

\begin{solution}{iq:qm:linear-models-under-stress:3}
La \omterm{def:qm:linear-models-under-stress:penalties}{régression ridge} rétrécit le long des petites directions singulières et répartit le poids sur un bloc corrélé~; le \omterm{def:qm:linear-models-under-stress:penalties}{lasso} annule des coefficients et tend à
conserver un seul élément d'un bloc, de façon quelque peu arbitraire~; l'\omterm{def:qm:linear-models-under-stress:penalties}{elastic net} garde les groupes ensemble tout en continuant à sélectionner.
\iqlookfor{Rétrécissement des coefficients contre sélection, et comportement au sein des groupes corrélés.}
\end{solution}

\begin{solution}{iq:qm:linear-models-under-stress:4}
Des plis dans l'ordre temporel, apprentissage avant test, avec un intervalle de purge d'au moins l'horizon de la cible (vingt jours) afin qu'il n'y ait aucune fuite entre cibles qui se chevauchent,
et éventuellement un embargo après le bloc de test~; standardiser et choisir les variables à l'intérieur de chaque pli d'apprentissage uniquement.
\iqlookfor{Ordre temporel, purge pour les chevauchements, et aucun prétraitement sur l'échantillon complet.}
\end{solution}

\begin{solution}{iq:qm:linear-models-under-stress:5}
Lorsque les variables explicatives et les résidus présentent des effets entreprise persistants~: la pente de chaque période partage la même erreur, si bien que l'écart type des séries temporelles
sous-estime l'incertitude. Il faut utiliser des \omterm{def:qm:linear-models-under-stress:fm}{erreurs types groupées} (clustered) par entreprise (ou par entreprise et par période lorsque les deux effets existent).
\iqlookfor{Quelle dépendance chaque méthode traite~: Fama--MacBeth traite les effets temporels, le groupement par entreprise traite les effets entreprise.}
\end{solution}

\begin{solution}{iq:qm:linear-models-under-stress:6}
Le coefficient de $X_2$ dans la régression sur $(X_1, X_2)$ égale la pente de $M_1y$ sur $M_1X_2$, soit les parties orthogonales à $X_1$~; preuve en appliquant
$M_1$ aux équations normales. Sur un desk~: le bêta d'une action par rapport à un secteur, une fois le marché retiré, est la pente des rendements résiduels au marché~; l'apport marginal d'un signal
après les signaux existants est son ajustement sur les rendements résiduels.
\iqlookfor{L'argument de projection et une orthogonalisation concrète.}
\end{solution}
